\section{Lean formalization and the remaining work}\label{sec:lean}

\subsection{The repository interface actually inspected}
The repository was inspected at commit
\texttt{0bba41d70f061442357cae0f6de055434a33fb09}.
Its root toolchain is Lean \texttt{v4.32.0}, and its Lake configuration
pins Mathlib \texttt{v4.32.0}. Relevant inspected modules include
\texttt{PR.lean}, \texttt{Syntax.lean}, \texttt{TeqReflect.lean}, and
\texttt{NativeBiInterpretation.lean}, together with the ConciseIPRA
README. The \texttt{docs/drafts} listing was inspected; none of its listed
archives was needed as an additional premise for these proofs.

The existing bridge is the following theorem, shown with ASCII
renderings of the repository notation:
\begin{lstlisting}
theorem pra_iff_native
  {Gamma : List (Fm PTm n)} {phi : Fm PTm n} :
  (Gamma |- [PRAAx] phi) <->
  NPrf (Gamma.map NFm.ofPRA) (NFm.ofPRA phi)
\end{lstlisting}
Its namespace is \texttt{LeanProofs.ConciseIPRA}. The theorem and formula
translations are in \texttt{NativeBiInterpretation.lean}
\cite{repo-bridge}. Thus an object-iPRA derivation, once constructed,
already has a route into the concise native calculus. A new proof of
bi-interpretability is unnecessary.

\subsection{What the supplied Lean sources contain}
The archive has two concrete source candidates for the ambient algebra:
\texttt{GeometryCore.lean} and \texttt{PascalConic.lean}. They use
explicit three- and six-component structures, dot and cross products,
a six-term $3\times3$ determinant, and a 20-term Laplace expansion of
the six-row determinant. The main equalities have \texttt{ring} proof
scripts; this is Mathlib's commutative ring normalization tactic
\cite{mathlib-ring}.

The source includes the full universal Desargues identity, its
center-to-collinearity implication and converse determinant implication,
the explicit-axis incidence lemma for an apart pair, the Pascal bracket
and quadratic-evaluation determinant identities, and the general-conic
Pascal implication. The last uses six column-replacement identities:
replacing column $j$ of $M$ by $Mq$ gives determinant $q_j\det M$.
The six conic equations make that replaced column zero; cancellation
then gives $\det M=0$. This is the same adjugate argument as the paper,
written without depending on a high-level matrix or geometry API.

\textbf{Validation status.} No Lean or Lake executable was available
in the execution environment. These new sources have \emph{not} been
compiled or kernel-checked. They contain no proof holes or added axiom
declarations, but that is not a substitute for compilation. Elaboration,
API, or performance repairs may still be needed. In particular, no new
object-theory theorem of the form $[]\vdash[\mathrm{PRAAx}]\varphi$ is
claimed as an already checked artifact. The written mathematics and
executed exact polynomial certificates have a separate status.

\subsection{A concrete path to a complete Lean proof}
The following milestones specify the remaining obligations.
\begin{enumerate}
\item \textbf{Compile the ambient algebra.} In a Lean/Mathlib
\texttt{v4.32.0} environment, compile the supplied files and repair any
elaboration issues. Audit the resulting theorems with
\texttt{\#print axioms}, rejecting \texttt{sorryAx}. This checks the
algebra but does not yet construct an iPRA derivation.

\item \textbf{Implement the numerical codes in the object language.}
Using \texttt{PR.zero}, \texttt{PR.succ}, \texttt{PR.proj},
\texttt{PR.comp}, and \texttt{PR.prec}, define the symbols for addition,
multiplication, monus, equality tests, signed pairs, fixed tuples, and
the four-branch axis function. The inspected proof interface provides
\texttt{PRA.proj}, \texttt{PRA.comp}, \texttt{PRA.precZero},
\texttt{PRA.precSucc}, and \texttt{PRA.qfInd}.

\item \textbf{Prove the arithmetic library internally.}
Construct \texttt{Prf PRAAx} derivations of the lemmas in
Section~\ref{sec:arith}, including signed-ring laws, decidability, and
no-zero-divisors. These are derivations with arbitrary object variables;
correctness only in the standard model would not suffice.

\item \textbf{Replay finite algebra in the object theory.}
Either implement a derivation-producing normalizer for finite polynomial
expressions or generate the fixed distributivity, associativity,
commutativity, and cancellation chains used here. A reflective route
must prove its normalizer sound for \texttt{Prf PRAAx}; an ambient
\texttt{ring} proof does not automatically translate itself. The exact
coefficient checks fix the identities and signs to be replayed.

\item \textbf{Construct the geometric derivations.}
Encode the formulas of Section~\ref{sec:internal} as \texttt{Fm PTm}.
Replay the incidence lemmas, finite coefficient selection, cancellation,
and four axis branches. The existing intuitionistic deduction
constructors handle universal closure and explicit existential witnesses.
This gives the actual target objects
\texttt{Prf PRAAx [] desarguesSentence} and
\texttt{Prf PRAAx [] pascalSentence}, where the two sentences are to
be defined by the encoding just described.

\item \textbf{Translate to the native calculus.}
Apply \texttt{pra\_iff\_native.mp} to those object-theory derivations.
The output is a native \texttt{NPrf} derivation of each translated sentence.

\item \textbf{Optional real schema.}
For a fixed PR coordinate family, encode the regularity, zero, and
apartness formulas and prove the finite threshold calculations in
Section~\ref{sec:real}. Reuse the same polynomial certificates. This
extension is separate from the exact rational theorem.
\end{enumerate}

The substantial remaining implementation is the object-arithmetic
library and replay of finite ring-equation derivations. There is no
unresolved geometric lemma in this plan. It would nevertheless be
inaccurate to call the ambient Lean candidate a completed formal proof
in the deeply embedded iPRA.

\section{Executable evidence and reproducibility}\label{sec:evidence}

The archive contains three complementary checks of the formulas.
\begin{enumerate}
\item The dependency-free sparse-polynomial verifier checks Pascal's
identity, its two-product bracket form, the four bracket relations, and
all six conic cofactor certificates. Variables are indeterminates;
coefficients are exact integers. The Pascal polynomial has exactly
$720$ nonzero monomials, with coefficients $+1$ or $-1$.

\item A second symbolic verifier checks the complete Desargues identity
and all three coordinates of the adjugate identity
\eqref{eq:adj3}. The Desargues polynomial has $1{,}290$ nonzero monomials.
A supplementary checker verifies the six generic column-replacement
identities used by the Lean candidate.

\item The integer geometry program checks the explicit examples and
the common-kernel algorithm on the cube $\{-1,0,1\}^3$.
Of $18{,}954$ input triples with $u\ne0$, exactly $7{,}146$ have zero
determinant; all returned vectors are nonzero and orthogonal to all
three inputs. The four branches are used $5{,}040$, $1{,}872$, $216$,
and $18$ times. This finite check exercises the implementation; the
universal proof is Lemma~\ref{lem:axis}.
\end{enumerate}
The logs and polynomial fingerprints are included. These computations
verify exact formulas; they are not emitted iPRA proofs or substitutes
for Lean kernel checking. The source code uses the Python standard
library only.

Run the checks using the commands in \texttt{README.txt}. Rebuild the
paper from \texttt{paper.tex} with \texttt{pdflatex} twice, or run
\texttt{build.sh}. The archive includes the uncompiled Lean sources,
pinned build metadata, a detailed formalization roadmap, and source
URLs with hashes for the repository audit.

\section{Conclusions}
Both theorems have proofs whose arithmetic core is finite
commutative-ring algebra. Desargues is governed by the product of the
two triangle determinants and the determinant of the vertex-joining
lines. Pascal is governed by the determinant of the six quadratic
evaluation rows, with a nonzero conic coefficient canceled through an
explicit cofactor certificate. These calculations give complete iPRA
proofs for integer-coded rational configurations, including primitive
recursive witnesses for their geometric conclusions. The real-name
schema makes the extra apartness information explicit. The Lean plan
preserves the distinction between checking the algebra in Lean and
constructing derivations inside intuitionistic primitive recursive
arithmetic.

